\documentclass[11pt]{article}

    \usepackage[breakable]{tcolorbox}
    \usepackage{parskip} % Stop auto-indenting (to mimic markdown behaviour)
    

    % Basic figure setup, for now with no caption control since it's done
    % automatically by Pandoc (which extracts ![](path) syntax from Markdown).
    \usepackage{graphicx}
    % Keep aspect ratio if custom image width or height is specified
    \setkeys{Gin}{keepaspectratio}
    % Maintain compatibility with old templates. Remove in nbconvert 6.0
    \let\Oldincludegraphics\includegraphics
    % Ensure that by default, figures have no caption (until we provide a
    % proper Figure object with a Caption API and a way to capture that
    % in the conversion process - todo).
    \usepackage{caption}
    \DeclareCaptionFormat{nocaption}{}
    \captionsetup{format=nocaption,aboveskip=0pt,belowskip=0pt}

    \usepackage{float}
    \floatplacement{figure}{H} % forces figures to be placed at the correct location
    \usepackage{xcolor} % Allow colors to be defined
    \usepackage{enumerate} % Needed for markdown enumerations to work
    \usepackage{geometry} % Used to adjust the document margins
    \usepackage{amsmath} % Equations
    \usepackage{amssymb} % Equations
    \usepackage{textcomp} % defines textquotesingle
    % Hack from http://tex.stackexchange.com/a/47451/13684:
    \AtBeginDocument{%
        \def\PYZsq{\textquotesingle}% Upright quotes in Pygmentized code
    }
    \usepackage{upquote} % Upright quotes for verbatim code
    \usepackage{eurosym} % defines \euro

    \usepackage{iftex}
    \ifPDFTeX
        \usepackage[T1]{fontenc}
        \IfFileExists{alphabeta.sty}{
              \usepackage{alphabeta}
          }{
              \usepackage[mathletters]{ucs}
              \usepackage[utf8x]{inputenc}
          }
    \else
        \usepackage{fontspec}
        \usepackage{unicode-math}
    \fi

    \usepackage{fancyvrb} % verbatim replacement that allows latex
    \usepackage{grffile} % extends the file name processing of package graphics
                         % to support a larger range
    \makeatletter % fix for old versions of grffile with XeLaTeX
    \@ifpackagelater{grffile}{2019/11/01}
    {
      % Do nothing on new versions
    }
    {
      \def\Gread@@xetex#1{%
        \IfFileExists{"\Gin@base".bb}%
        {\Gread@eps{\Gin@base.bb}}%
        {\Gread@@xetex@aux#1}%
      }
    }
    \makeatother
    \usepackage[Export]{adjustbox} % Used to constrain images to a maximum size
    \adjustboxset{max size={0.9\linewidth}{0.9\paperheight}}

    % The hyperref package gives us a pdf with properly built
    % internal navigation ('pdf bookmarks' for the table of contents,
    % internal cross-reference links, web links for URLs, etc.)
    \usepackage{hyperref}
    % The default LaTeX title has an obnoxious amount of whitespace. By default,
    % titling removes some of it. It also provides customization options.
    \usepackage{titling}
    \usepackage{longtable} % longtable support required by pandoc >1.10
    \usepackage{booktabs}  % table support for pandoc > 1.12.2
    \usepackage{array}     % table support for pandoc >= 2.11.3
    \usepackage{calc}      % table minipage width calculation for pandoc >= 2.11.1
    \usepackage[inline]{enumitem} % IRkernel/repr support (it uses the enumerate* environment)
    \usepackage[normalem]{ulem} % ulem is needed to support strikethroughs (\sout)
                                % normalem makes italics be italics, not underlines
    \usepackage{soul}      % strikethrough (\st) support for pandoc >= 3.0.0
    \usepackage{mathrsfs}
    

    
    % Colors for the hyperref package
    \definecolor{urlcolor}{rgb}{0,.145,.698}
    \definecolor{linkcolor}{rgb}{.71,0.21,0.01}
    \definecolor{citecolor}{rgb}{.12,.54,.11}

    % ANSI colors
    \definecolor{ansi-black}{HTML}{3E424D}
    \definecolor{ansi-black-intense}{HTML}{282C36}
    \definecolor{ansi-red}{HTML}{E75C58}
    \definecolor{ansi-red-intense}{HTML}{B22B31}
    \definecolor{ansi-green}{HTML}{00A250}
    \definecolor{ansi-green-intense}{HTML}{007427}
    \definecolor{ansi-yellow}{HTML}{DDB62B}
    \definecolor{ansi-yellow-intense}{HTML}{B27D12}
    \definecolor{ansi-blue}{HTML}{208FFB}
    \definecolor{ansi-blue-intense}{HTML}{0065CA}
    \definecolor{ansi-magenta}{HTML}{D160C4}
    \definecolor{ansi-magenta-intense}{HTML}{A03196}
    \definecolor{ansi-cyan}{HTML}{60C6C8}
    \definecolor{ansi-cyan-intense}{HTML}{258F8F}
    \definecolor{ansi-white}{HTML}{C5C1B4}
    \definecolor{ansi-white-intense}{HTML}{A1A6B2}
    \definecolor{ansi-default-inverse-fg}{HTML}{FFFFFF}
    \definecolor{ansi-default-inverse-bg}{HTML}{000000}

    % common color for the border for error outputs.
    \definecolor{outerrorbackground}{HTML}{FFDFDF}

    % commands and environments needed by pandoc snippets
    % extracted from the output of `pandoc -s`
    \providecommand{\tightlist}{%
      \setlength{\itemsep}{0pt}\setlength{\parskip}{0pt}}
    \DefineVerbatimEnvironment{Highlighting}{Verbatim}{commandchars=\\\{\}}
    % Add ',fontsize=\small' for more characters per line
    \newenvironment{Shaded}{}{}
    \newcommand{\KeywordTok}[1]{\textcolor[rgb]{0.00,0.44,0.13}{\textbf{{#1}}}}
    \newcommand{\DataTypeTok}[1]{\textcolor[rgb]{0.56,0.13,0.00}{{#1}}}
    \newcommand{\DecValTok}[1]{\textcolor[rgb]{0.25,0.63,0.44}{{#1}}}
    \newcommand{\BaseNTok}[1]{\textcolor[rgb]{0.25,0.63,0.44}{{#1}}}
    \newcommand{\FloatTok}[1]{\textcolor[rgb]{0.25,0.63,0.44}{{#1}}}
    \newcommand{\CharTok}[1]{\textcolor[rgb]{0.25,0.44,0.63}{{#1}}}
    \newcommand{\StringTok}[1]{\textcolor[rgb]{0.25,0.44,0.63}{{#1}}}
    \newcommand{\CommentTok}[1]{\textcolor[rgb]{0.38,0.63,0.69}{\textit{{#1}}}}
    \newcommand{\OtherTok}[1]{\textcolor[rgb]{0.00,0.44,0.13}{{#1}}}
    \newcommand{\AlertTok}[1]{\textcolor[rgb]{1.00,0.00,0.00}{\textbf{{#1}}}}
    \newcommand{\FunctionTok}[1]{\textcolor[rgb]{0.02,0.16,0.49}{{#1}}}
    \newcommand{\RegionMarkerTok}[1]{{#1}}
    \newcommand{\ErrorTok}[1]{\textcolor[rgb]{1.00,0.00,0.00}{\textbf{{#1}}}}
    \newcommand{\NormalTok}[1]{{#1}}

    % Additional commands for more recent versions of Pandoc
    \newcommand{\ConstantTok}[1]{\textcolor[rgb]{0.53,0.00,0.00}{{#1}}}
    \newcommand{\SpecialCharTok}[1]{\textcolor[rgb]{0.25,0.44,0.63}{{#1}}}
    \newcommand{\VerbatimStringTok}[1]{\textcolor[rgb]{0.25,0.44,0.63}{{#1}}}
    \newcommand{\SpecialStringTok}[1]{\textcolor[rgb]{0.73,0.40,0.53}{{#1}}}
    \newcommand{\ImportTok}[1]{{#1}}
    \newcommand{\DocumentationTok}[1]{\textcolor[rgb]{0.73,0.13,0.13}{\textit{{#1}}}}
    \newcommand{\AnnotationTok}[1]{\textcolor[rgb]{0.38,0.63,0.69}{\textbf{\textit{{#1}}}}}
    \newcommand{\CommentVarTok}[1]{\textcolor[rgb]{0.38,0.63,0.69}{\textbf{\textit{{#1}}}}}
    \newcommand{\VariableTok}[1]{\textcolor[rgb]{0.10,0.09,0.49}{{#1}}}
    \newcommand{\ControlFlowTok}[1]{\textcolor[rgb]{0.00,0.44,0.13}{\textbf{{#1}}}}
    \newcommand{\OperatorTok}[1]{\textcolor[rgb]{0.40,0.40,0.40}{{#1}}}
    \newcommand{\BuiltInTok}[1]{{#1}}
    \newcommand{\ExtensionTok}[1]{{#1}}
    \newcommand{\PreprocessorTok}[1]{\textcolor[rgb]{0.74,0.48,0.00}{{#1}}}
    \newcommand{\AttributeTok}[1]{\textcolor[rgb]{0.49,0.56,0.16}{{#1}}}
    \newcommand{\InformationTok}[1]{\textcolor[rgb]{0.38,0.63,0.69}{\textbf{\textit{{#1}}}}}
    \newcommand{\WarningTok}[1]{\textcolor[rgb]{0.38,0.63,0.69}{\textbf{\textit{{#1}}}}}


    % Define a nice break command that doesn't care if a line doesn't already
    % exist.
    \def\br{\hspace*{\fill} \\* }
    % Math Jax compatibility definitions
    \def\gt{>}
    \def\lt{<}
    \let\Oldtex\TeX
    \let\Oldlatex\LaTeX
    \renewcommand{\TeX}{\textrm{\Oldtex}}
    \renewcommand{\LaTeX}{\textrm{\Oldlatex}}
    % Document parameters
    % Document title
    \title{3b. Overall\_exam\_performance\_analysis\_genders}
    
    
    
    
    
    
    
% Pygments definitions
\makeatletter
\def\PY@reset{\let\PY@it=\relax \let\PY@bf=\relax%
    \let\PY@ul=\relax \let\PY@tc=\relax%
    \let\PY@bc=\relax \let\PY@ff=\relax}
\def\PY@tok#1{\csname PY@tok@#1\endcsname}
\def\PY@toks#1+{\ifx\relax#1\empty\else%
    \PY@tok{#1}\expandafter\PY@toks\fi}
\def\PY@do#1{\PY@bc{\PY@tc{\PY@ul{%
    \PY@it{\PY@bf{\PY@ff{#1}}}}}}}
\def\PY#1#2{\PY@reset\PY@toks#1+\relax+\PY@do{#2}}

\@namedef{PY@tok@w}{\def\PY@tc##1{\textcolor[rgb]{0.73,0.73,0.73}{##1}}}
\@namedef{PY@tok@c}{\let\PY@it=\textit\def\PY@tc##1{\textcolor[rgb]{0.24,0.48,0.48}{##1}}}
\@namedef{PY@tok@cp}{\def\PY@tc##1{\textcolor[rgb]{0.61,0.40,0.00}{##1}}}
\@namedef{PY@tok@k}{\let\PY@bf=\textbf\def\PY@tc##1{\textcolor[rgb]{0.00,0.50,0.00}{##1}}}
\@namedef{PY@tok@kp}{\def\PY@tc##1{\textcolor[rgb]{0.00,0.50,0.00}{##1}}}
\@namedef{PY@tok@kt}{\def\PY@tc##1{\textcolor[rgb]{0.69,0.00,0.25}{##1}}}
\@namedef{PY@tok@o}{\def\PY@tc##1{\textcolor[rgb]{0.40,0.40,0.40}{##1}}}
\@namedef{PY@tok@ow}{\let\PY@bf=\textbf\def\PY@tc##1{\textcolor[rgb]{0.67,0.13,1.00}{##1}}}
\@namedef{PY@tok@nb}{\def\PY@tc##1{\textcolor[rgb]{0.00,0.50,0.00}{##1}}}
\@namedef{PY@tok@nf}{\def\PY@tc##1{\textcolor[rgb]{0.00,0.00,1.00}{##1}}}
\@namedef{PY@tok@nc}{\let\PY@bf=\textbf\def\PY@tc##1{\textcolor[rgb]{0.00,0.00,1.00}{##1}}}
\@namedef{PY@tok@nn}{\let\PY@bf=\textbf\def\PY@tc##1{\textcolor[rgb]{0.00,0.00,1.00}{##1}}}
\@namedef{PY@tok@ne}{\let\PY@bf=\textbf\def\PY@tc##1{\textcolor[rgb]{0.80,0.25,0.22}{##1}}}
\@namedef{PY@tok@nv}{\def\PY@tc##1{\textcolor[rgb]{0.10,0.09,0.49}{##1}}}
\@namedef{PY@tok@no}{\def\PY@tc##1{\textcolor[rgb]{0.53,0.00,0.00}{##1}}}
\@namedef{PY@tok@nl}{\def\PY@tc##1{\textcolor[rgb]{0.46,0.46,0.00}{##1}}}
\@namedef{PY@tok@ni}{\let\PY@bf=\textbf\def\PY@tc##1{\textcolor[rgb]{0.44,0.44,0.44}{##1}}}
\@namedef{PY@tok@na}{\def\PY@tc##1{\textcolor[rgb]{0.41,0.47,0.13}{##1}}}
\@namedef{PY@tok@nt}{\let\PY@bf=\textbf\def\PY@tc##1{\textcolor[rgb]{0.00,0.50,0.00}{##1}}}
\@namedef{PY@tok@nd}{\def\PY@tc##1{\textcolor[rgb]{0.67,0.13,1.00}{##1}}}
\@namedef{PY@tok@s}{\def\PY@tc##1{\textcolor[rgb]{0.73,0.13,0.13}{##1}}}
\@namedef{PY@tok@sd}{\let\PY@it=\textit\def\PY@tc##1{\textcolor[rgb]{0.73,0.13,0.13}{##1}}}
\@namedef{PY@tok@si}{\let\PY@bf=\textbf\def\PY@tc##1{\textcolor[rgb]{0.64,0.35,0.47}{##1}}}
\@namedef{PY@tok@se}{\let\PY@bf=\textbf\def\PY@tc##1{\textcolor[rgb]{0.67,0.36,0.12}{##1}}}
\@namedef{PY@tok@sr}{\def\PY@tc##1{\textcolor[rgb]{0.64,0.35,0.47}{##1}}}
\@namedef{PY@tok@ss}{\def\PY@tc##1{\textcolor[rgb]{0.10,0.09,0.49}{##1}}}
\@namedef{PY@tok@sx}{\def\PY@tc##1{\textcolor[rgb]{0.00,0.50,0.00}{##1}}}
\@namedef{PY@tok@m}{\def\PY@tc##1{\textcolor[rgb]{0.40,0.40,0.40}{##1}}}
\@namedef{PY@tok@gh}{\let\PY@bf=\textbf\def\PY@tc##1{\textcolor[rgb]{0.00,0.00,0.50}{##1}}}
\@namedef{PY@tok@gu}{\let\PY@bf=\textbf\def\PY@tc##1{\textcolor[rgb]{0.50,0.00,0.50}{##1}}}
\@namedef{PY@tok@gd}{\def\PY@tc##1{\textcolor[rgb]{0.63,0.00,0.00}{##1}}}
\@namedef{PY@tok@gi}{\def\PY@tc##1{\textcolor[rgb]{0.00,0.52,0.00}{##1}}}
\@namedef{PY@tok@gr}{\def\PY@tc##1{\textcolor[rgb]{0.89,0.00,0.00}{##1}}}
\@namedef{PY@tok@ge}{\let\PY@it=\textit}
\@namedef{PY@tok@gs}{\let\PY@bf=\textbf}
\@namedef{PY@tok@gp}{\let\PY@bf=\textbf\def\PY@tc##1{\textcolor[rgb]{0.00,0.00,0.50}{##1}}}
\@namedef{PY@tok@go}{\def\PY@tc##1{\textcolor[rgb]{0.44,0.44,0.44}{##1}}}
\@namedef{PY@tok@gt}{\def\PY@tc##1{\textcolor[rgb]{0.00,0.27,0.87}{##1}}}
\@namedef{PY@tok@err}{\def\PY@bc##1{{\setlength{\fboxsep}{\string -\fboxrule}\fcolorbox[rgb]{1.00,0.00,0.00}{1,1,1}{\strut ##1}}}}
\@namedef{PY@tok@kc}{\let\PY@bf=\textbf\def\PY@tc##1{\textcolor[rgb]{0.00,0.50,0.00}{##1}}}
\@namedef{PY@tok@kd}{\let\PY@bf=\textbf\def\PY@tc##1{\textcolor[rgb]{0.00,0.50,0.00}{##1}}}
\@namedef{PY@tok@kn}{\let\PY@bf=\textbf\def\PY@tc##1{\textcolor[rgb]{0.00,0.50,0.00}{##1}}}
\@namedef{PY@tok@kr}{\let\PY@bf=\textbf\def\PY@tc##1{\textcolor[rgb]{0.00,0.50,0.00}{##1}}}
\@namedef{PY@tok@bp}{\def\PY@tc##1{\textcolor[rgb]{0.00,0.50,0.00}{##1}}}
\@namedef{PY@tok@fm}{\def\PY@tc##1{\textcolor[rgb]{0.00,0.00,1.00}{##1}}}
\@namedef{PY@tok@vc}{\def\PY@tc##1{\textcolor[rgb]{0.10,0.09,0.49}{##1}}}
\@namedef{PY@tok@vg}{\def\PY@tc##1{\textcolor[rgb]{0.10,0.09,0.49}{##1}}}
\@namedef{PY@tok@vi}{\def\PY@tc##1{\textcolor[rgb]{0.10,0.09,0.49}{##1}}}
\@namedef{PY@tok@vm}{\def\PY@tc##1{\textcolor[rgb]{0.10,0.09,0.49}{##1}}}
\@namedef{PY@tok@sa}{\def\PY@tc##1{\textcolor[rgb]{0.73,0.13,0.13}{##1}}}
\@namedef{PY@tok@sb}{\def\PY@tc##1{\textcolor[rgb]{0.73,0.13,0.13}{##1}}}
\@namedef{PY@tok@sc}{\def\PY@tc##1{\textcolor[rgb]{0.73,0.13,0.13}{##1}}}
\@namedef{PY@tok@dl}{\def\PY@tc##1{\textcolor[rgb]{0.73,0.13,0.13}{##1}}}
\@namedef{PY@tok@s2}{\def\PY@tc##1{\textcolor[rgb]{0.73,0.13,0.13}{##1}}}
\@namedef{PY@tok@sh}{\def\PY@tc##1{\textcolor[rgb]{0.73,0.13,0.13}{##1}}}
\@namedef{PY@tok@s1}{\def\PY@tc##1{\textcolor[rgb]{0.73,0.13,0.13}{##1}}}
\@namedef{PY@tok@mb}{\def\PY@tc##1{\textcolor[rgb]{0.40,0.40,0.40}{##1}}}
\@namedef{PY@tok@mf}{\def\PY@tc##1{\textcolor[rgb]{0.40,0.40,0.40}{##1}}}
\@namedef{PY@tok@mh}{\def\PY@tc##1{\textcolor[rgb]{0.40,0.40,0.40}{##1}}}
\@namedef{PY@tok@mi}{\def\PY@tc##1{\textcolor[rgb]{0.40,0.40,0.40}{##1}}}
\@namedef{PY@tok@il}{\def\PY@tc##1{\textcolor[rgb]{0.40,0.40,0.40}{##1}}}
\@namedef{PY@tok@mo}{\def\PY@tc##1{\textcolor[rgb]{0.40,0.40,0.40}{##1}}}
\@namedef{PY@tok@ch}{\let\PY@it=\textit\def\PY@tc##1{\textcolor[rgb]{0.24,0.48,0.48}{##1}}}
\@namedef{PY@tok@cm}{\let\PY@it=\textit\def\PY@tc##1{\textcolor[rgb]{0.24,0.48,0.48}{##1}}}
\@namedef{PY@tok@cpf}{\let\PY@it=\textit\def\PY@tc##1{\textcolor[rgb]{0.24,0.48,0.48}{##1}}}
\@namedef{PY@tok@c1}{\let\PY@it=\textit\def\PY@tc##1{\textcolor[rgb]{0.24,0.48,0.48}{##1}}}
\@namedef{PY@tok@cs}{\let\PY@it=\textit\def\PY@tc##1{\textcolor[rgb]{0.24,0.48,0.48}{##1}}}

\def\PYZbs{\char`\\}
\def\PYZus{\char`\_}
\def\PYZob{\char`\{}
\def\PYZcb{\char`\}}
\def\PYZca{\char`\^}
\def\PYZam{\char`\&}
\def\PYZlt{\char`\<}
\def\PYZgt{\char`\>}
\def\PYZsh{\char`\#}
\def\PYZpc{\char`\%}
\def\PYZdl{\char`\$}
\def\PYZhy{\char`\-}
\def\PYZsq{\char`\'}
\def\PYZdq{\char`\"}
\def\PYZti{\char`\~}
% for compatibility with earlier versions
\def\PYZat{@}
\def\PYZlb{[}
\def\PYZrb{]}
\makeatother


    % For linebreaks inside Verbatim environment from package fancyvrb.
    \makeatletter
        \newbox\Wrappedcontinuationbox
        \newbox\Wrappedvisiblespacebox
        \newcommand*\Wrappedvisiblespace {\textcolor{red}{\textvisiblespace}}
        \newcommand*\Wrappedcontinuationsymbol {\textcolor{red}{\llap{\tiny$\m@th\hookrightarrow$}}}
        \newcommand*\Wrappedcontinuationindent {3ex }
        \newcommand*\Wrappedafterbreak {\kern\Wrappedcontinuationindent\copy\Wrappedcontinuationbox}
        % Take advantage of the already applied Pygments mark-up to insert
        % potential linebreaks for TeX processing.
        %        {, <, #, %, $, ' and ": go to next line.
        %        _, }, ^, &, >, - and ~: stay at end of broken line.
        % Use of \textquotesingle for straight quote.
        \newcommand*\Wrappedbreaksatspecials {%
            \def\PYGZus{\discretionary{\char`\_}{\Wrappedafterbreak}{\char`\_}}%
            \def\PYGZob{\discretionary{}{\Wrappedafterbreak\char`\{}{\char`\{}}%
            \def\PYGZcb{\discretionary{\char`\}}{\Wrappedafterbreak}{\char`\}}}%
            \def\PYGZca{\discretionary{\char`\^}{\Wrappedafterbreak}{\char`\^}}%
            \def\PYGZam{\discretionary{\char`\&}{\Wrappedafterbreak}{\char`\&}}%
            \def\PYGZlt{\discretionary{}{\Wrappedafterbreak\char`\<}{\char`\<}}%
            \def\PYGZgt{\discretionary{\char`\>}{\Wrappedafterbreak}{\char`\>}}%
            \def\PYGZsh{\discretionary{}{\Wrappedafterbreak\char`\#}{\char`\#}}%
            \def\PYGZpc{\discretionary{}{\Wrappedafterbreak\char`\%}{\char`\%}}%
            \def\PYGZdl{\discretionary{}{\Wrappedafterbreak\char`\$}{\char`\$}}%
            \def\PYGZhy{\discretionary{\char`\-}{\Wrappedafterbreak}{\char`\-}}%
            \def\PYGZsq{\discretionary{}{\Wrappedafterbreak\textquotesingle}{\textquotesingle}}%
            \def\PYGZdq{\discretionary{}{\Wrappedafterbreak\char`\"}{\char`\"}}%
            \def\PYGZti{\discretionary{\char`\~}{\Wrappedafterbreak}{\char`\~}}%
        }
        % Some characters . , ; ? ! / are not pygmentized.
        % This macro makes them "active" and they will insert potential linebreaks
        \newcommand*\Wrappedbreaksatpunct {%
            \lccode`\~`\.\lowercase{\def~}{\discretionary{\hbox{\char`\.}}{\Wrappedafterbreak}{\hbox{\char`\.}}}%
            \lccode`\~`\,\lowercase{\def~}{\discretionary{\hbox{\char`\,}}{\Wrappedafterbreak}{\hbox{\char`\,}}}%
            \lccode`\~`\;\lowercase{\def~}{\discretionary{\hbox{\char`\;}}{\Wrappedafterbreak}{\hbox{\char`\;}}}%
            \lccode`\~`\:\lowercase{\def~}{\discretionary{\hbox{\char`\:}}{\Wrappedafterbreak}{\hbox{\char`\:}}}%
            \lccode`\~`\?\lowercase{\def~}{\discretionary{\hbox{\char`\?}}{\Wrappedafterbreak}{\hbox{\char`\?}}}%
            \lccode`\~`\!\lowercase{\def~}{\discretionary{\hbox{\char`\!}}{\Wrappedafterbreak}{\hbox{\char`\!}}}%
            \lccode`\~`\/\lowercase{\def~}{\discretionary{\hbox{\char`\/}}{\Wrappedafterbreak}{\hbox{\char`\/}}}%
            \catcode`\.\active
            \catcode`\,\active
            \catcode`\;\active
            \catcode`\:\active
            \catcode`\?\active
            \catcode`\!\active
            \catcode`\/\active
            \lccode`\~`\~
        }
    \makeatother

    \let\OriginalVerbatim=\Verbatim
    \makeatletter
    \renewcommand{\Verbatim}[1][1]{%
        %\parskip\z@skip
        \sbox\Wrappedcontinuationbox {\Wrappedcontinuationsymbol}%
        \sbox\Wrappedvisiblespacebox {\FV@SetupFont\Wrappedvisiblespace}%
        \def\FancyVerbFormatLine ##1{\hsize\linewidth
            \vtop{\raggedright\hyphenpenalty\z@\exhyphenpenalty\z@
                \doublehyphendemerits\z@\finalhyphendemerits\z@
                \strut ##1\strut}%
        }%
        % If the linebreak is at a space, the latter will be displayed as visible
        % space at end of first line, and a continuation symbol starts next line.
        % Stretch/shrink are however usually zero for typewriter font.
        \def\FV@Space {%
            \nobreak\hskip\z@ plus\fontdimen3\font minus\fontdimen4\font
            \discretionary{\copy\Wrappedvisiblespacebox}{\Wrappedafterbreak}
            {\kern\fontdimen2\font}%
        }%

        % Allow breaks at special characters using \PYG... macros.
        \Wrappedbreaksatspecials
        % Breaks at punctuation characters . , ; ? ! and / need catcode=\active
        \OriginalVerbatim[#1,codes*=\Wrappedbreaksatpunct]%
    }
    \makeatother

    % Exact colors from NB
    \definecolor{incolor}{HTML}{303F9F}
    \definecolor{outcolor}{HTML}{D84315}
    \definecolor{cellborder}{HTML}{CFCFCF}
    \definecolor{cellbackground}{HTML}{F7F7F7}

    % prompt
    \makeatletter
    \newcommand{\boxspacing}{\kern\kvtcb@left@rule\kern\kvtcb@boxsep}
    \makeatother
    \newcommand{\prompt}[4]{
        {\ttfamily\llap{{\color{#2}[#3]:\hspace{3pt}#4}}\vspace{-\baselineskip}}
    }
    

    
    % Prevent overflowing lines due to hard-to-break entities
    \sloppy
    % Setup hyperref package
    \hypersetup{
      breaklinks=true,  % so long urls are correctly broken across lines
      colorlinks=true,
      urlcolor=urlcolor,
      linkcolor=linkcolor,
      citecolor=citecolor,
      }
    % Slightly bigger margins than the latex defaults
    
    \geometry{verbose,tmargin=1in,bmargin=1in,lmargin=1in,rmargin=1in}
    
    

\begin{document}
    
    \maketitle
    
    

    
    \section{Exam Performance Analysis: Gender and Time
Patterns}\label{exam-performance-analysis-gender-and-time-patterns}

This notebook analyzes exam performance differences between male and
female students, focusing on:

\begin{itemize}
\tightlist
\item
  \textbf{Performance differences} between genders with statistical
  testing
\item
  \textbf{Time completion patterns} by gender
\item
  \textbf{Effect size} and practical significance of observed
  differences
\item
  \textbf{Time-performance correlations} and their relationship to
  gender
\end{itemize}

    \subsection{1. Data Loading}\label{data-loading}

The cleaned exam data is loaded from the CSV file.

    \begin{Verbatim}[commandchars=\\\{\}]
Data loaded successfully from ../data/cleaned\_data/IN1020\_2024\_cleaned\_data.csv
    \end{Verbatim}

    \subsection{Overview}\label{overview}

    \begin{Verbatim}[commandchars=\\\{\}]
This notebook provides a comprehensive analysis of exam performance for
IN1020\_2024, comparing male and female students.
    \end{Verbatim}

    \subsection{2. Data Preparation}\label{data-preparation}

We'll prepare separate datasets for different analyses: -
\textbf{df\_male}: Data about total score among male students -
\textbf{df\_female}: Data about total score among female students -
\textbf{df\_time}: Data including exam timing information by gender

    \begin{Verbatim}[commandchars=\\\{\}]
Dataset with only male students:
    \end{Verbatim}

            \begin{tcolorbox}[breakable, size=fbox, boxrule=.5pt, pad at break*=1mm, opacityfill=0]
\begin{Verbatim}[commandchars=\\\{\}]
   student\_id gender  total\_score
0           1      M        72.63
2           3      M        75.45
4           5      M        72.28
6           7      M        56.61
8           9      M        69.36
\end{Verbatim}
\end{tcolorbox}
        
    \begin{Verbatim}[commandchars=\\\{\}]
Dataset with only female students:
    \end{Verbatim}

            \begin{tcolorbox}[breakable, size=fbox, boxrule=.5pt, pad at break*=1mm, opacityfill=0]
\begin{Verbatim}[commandchars=\\\{\}]
    student\_id gender  total\_score
1            2      F        56.71
3            4      F        55.88
5            6      F        73.43
7            8      F        61.29
10          11      F        48.38
\end{Verbatim}
\end{tcolorbox}
        
    \begin{Verbatim}[commandchars=\\\{\}]
Dataset with time taken to complete the exam:
    \end{Verbatim}

            \begin{tcolorbox}[breakable, size=fbox, boxrule=.5pt, pad at break*=1mm, opacityfill=0]
\begin{Verbatim}[commandchars=\\\{\}]
   student\_id      exam\_start\_time        exam\_end\_time gender  total\_score  \textbackslash{}
0           1  2024-12-11 14:01:44  2024-12-11 16:18:42      M        72.63
1           2  2024-12-11 14:00:02  2024-12-11 17:52:41      F        56.71
2           3  2024-12-11 14:00:03  2024-12-11 17:57:49      M        75.45
3           4  2024-12-11 14:00:05  2024-12-11 17:53:27      F        55.88
4           5  2024-12-11 14:00:04  2024-12-11 17:23:20      M        72.28

   extra\_time\_mins  incident\_time\_mins  time\_taken
0                0                   0      136.97
1                0                   0      232.65
2                0                   0      237.77
3                0                   0      233.37
4                0                   0      203.27
\end{Verbatim}
\end{tcolorbox}
        
    \begin{center}\rule{0.5\linewidth}{0.5pt}\end{center}

    \subsection{3. Descriptive Statistics}\label{descriptive-statistics}

This section presents comprehensive statistical measures for total
scores by gender, including: - Central tendency measures (mean, median,
mode) - Dispersion measures (standard deviation, range, IQR) -
Distribution shape indicators (skewness, kurtosis) - Gender performance
comparison

    \begin{Verbatim}[commandchars=\\\{\}]
=== DESCRIPTIVE STATISTICS FOR TOTAL SCORES ===
             Male   Female
Count     311.000  258.000
Mean       66.070   63.440
Median     67.930   66.660
Mode       63.530   42.500
Std Dev    13.330   14.300
Min        27.700   24.480
Max        91.850   90.890
Q1         58.140   53.810
Q3         75.620   74.140
IQR        17.490   20.330
Skewness   -0.656   -0.465
Kurtosis    0.005   -0.308

=== KEY FINDINGS ===
Mean difference (Male - Female): 2.63 points
Percentage difference: 4.1\%
Male performance is higher than female performance
    \end{Verbatim}

    \subsubsection{3.1 Gender Distribution}\label{gender-distribution}

The following pie chart shows the proportion of male and female students
in the exam.

    \begin{center}
    \adjustimage{max size={0.9\linewidth}{0.9\paperheight}}{3b. Overall_exam_performance_analysis_genders_files/3b. Overall_exam_performance_analysis_genders_16_0.png}
    \end{center}
    { \hspace*{\fill} \\}
    
    \begin{center}\rule{0.5\linewidth}{0.5pt}\end{center}

\subsection{4. Statistical Analysis}\label{statistical-analysis}

This section performs statistical hypothesis testing to determine if
performance differences between genders are statistically significant
and practically meaningful.

    \subsubsection{4.1 Score Distribution
Visualization}\label{score-distribution-visualization}

This section visualizes the score distributions for male and female
students with key statistical measures overlaid: - \textbf{Red dashed
line}: Mean score - \textbf{Green solid line}: Median score -
\textbf{Purple shaded area}: ±1 standard deviation from the mean -
\textbf{Skewness value}: Indicates distribution shape

    \begin{center}
    \adjustimage{max size={0.9\linewidth}{0.9\paperheight}}{3b. Overall_exam_performance_analysis_genders_files/3b. Overall_exam_performance_analysis_genders_19_0.png}
    \end{center}
    { \hspace*{\fill} \\}
    
    \begin{Verbatim}[commandchars=\\\{\}]
=== DISTRIBUTION SHAPE INTERPRETATION ===
Male: left-skewed (negative skew) (skewness = -0.656)
Female: approximately symmetric (skewness = -0.465)
    \end{Verbatim}

    \subsubsection{4.2 Data Distribution \& Normality
Testing}\label{data-distribution-normality-testing}

Before choosing the appropriate statistical test, we need to check if
the data follows a normal distribution.

    \begin{Verbatim}[commandchars=\\\{\}]
=== NORMALITY TESTS ===

Male (n=311):
  Shapiro-Wilk test for samples between 200 and 2000 entities: W=0.9646,
p=0.0000
  Normal distribution: False

Female (n=258):
  Shapiro-Wilk test for samples between 200 and 2000 entities: W=0.9762,
p=0.0003
  Normal distribution: False

Both groups normally distributed: False
    \end{Verbatim}

    \subsubsection{4.3 Hypothesis Testing}\label{hypothesis-testing}

The following analysis tests whether the observed performance difference
between male and female students is statistically significant and
calculates effect sizes to determine practical importance.

    \begin{Verbatim}[commandchars=\\\{\}]
=== STATISTICAL HYPOTHESIS TESTING ===
H0: No difference in mean scores between male and female students
H1: There is a difference in mean scores between groups

Primary approach: Non-parametric (based on normality test results)

NON-PARAMETRIC APPROACH (Data violates normality assumptions):
  Mann-Whitney U test: U = 44500.5, p = 0.0248
  Rank-biserial correlation: 0.1092
  Effect size magnitude: small
  Significant at α=0.05: True

CORRELATION ANALYSIS:
  Point-biserial correlation: r = 0.0948, p = 0.0237
  Variance explained (R²): 0.0090 (0.9\%)

PRIMARY INTERPRETATION:
  Based on Mann-Whitney U test:
  Statistical significance: Yes (p = 0.0248)
  Effect size: Rank-biserial r = 0.109 (small)
  Practical significance: Yes
  Decision: Reject H0 - There is a statistically significant difference between
groups
    \end{Verbatim}

    \subsubsection{4.4 Results Interpretation}\label{results-interpretation}

The following interpretation explains the statistical findings in
accessible language:

    \begin{Verbatim}[commandchars=\\\{\}]
UNDERSTANDING YOUR TEST RESULTS
======================================================================

 WHAT WE FOUND:
   • Male students averaged: 66.1 points
   • Female students averaged: 63.4 points
   • Difference: 2.6 points (male students scored higher)

 QUESTION 1: Is this difference REAL or just random chance?
   ANSWER: We are confident (95\%+ certain)
   What this means: This is likely a real difference, not random variation.
   Technical note: p-value = 0.0248

 QUESTION 2: How BIG or IMPORTANT is this difference?
   ANSWER: SMALL - noticeable but minor
   Suggested action: Worth monitoring but probably doesn't require immediate
intervention.

QUESTION 3: What does this mean in PRACTICAL terms?
   • The male group scores about 4.1\% higher on average
   • Individual students can still perform very differently from their group
average

SUMMARY:
   STATISTICALLY SIGNIFICANT BUT PRACTICALLY SMALL
   There IS a real difference, but it's so small it may not matter
educationally.
   Recommendation: Monitor future data but don't overreact to this small
difference.

======================================================================
    \end{Verbatim}

    \subsubsection{4.5 Distribution Visualization \& Outlier
Detection}\label{distribution-visualization-outlier-detection}

The following visualizations help identify outliers and understand the
distribution shape through boxplots and violin plots.

    \begin{Verbatim}[commandchars=\\\{\}]
=== OUTLIER ANALYSIS ===
Male: 4 outliers (1.3\%)
  Range: 27.70 - 30.02
Female: 0 outliers (0.0\%)
    \end{Verbatim}

    \begin{center}
    \adjustimage{max size={0.9\linewidth}{0.9\paperheight}}{3b. Overall_exam_performance_analysis_genders_files/3b. Overall_exam_performance_analysis_genders_27_1.png}
    \end{center}
    { \hspace*{\fill} \\}
    
    \begin{center}\rule{0.5\linewidth}{0.5pt}\end{center}

\subsection{5. Time-Performance Correlation
Analysis}\label{time-performance-correlation-analysis}

This section examines the relationship between time spent on the exam
and performance, analyzing: - Overall correlation between time and
scores - Gender-specific time-score relationships - Statistical
significance and practical implications

    \subsubsection{5.1 Correlation Analysis}\label{correlation-analysis}

    \begin{Verbatim}[commandchars=\\\{\}]
=== TIME-SCORE CORRELATION ANALYSIS ===

Testing normality for Time taken and Total score:

Time taken (n=569):
  Shapiro-Wilk test for samples between 200 and 2000 entities: W=0.9080,
p=0.0000
  Normal distribution: False

Total score (n=569):
  Shapiro-Wilk test for samples between 200 and 2000 entities: W=0.9709,
p=0.0000
  Normal distribution: False

Both variables normally distributed: False

Using NON-PARAMETRIC approach (Spearman correlation):
  Spearman correlation: ρ = -0.1349, p = 0.0013
  Pearson correlation (supplementary): r = -0.0839, p = 0.0456

PRIMARY INTERPRETATION (Spearman correlation):
  Correlation strength: weak negative correlation
  Statistical significance: Yes (p = 0.0013)

=== GENDER-SPECIFIC CORRELATIONS ===

Testing normality for Male time taken and Male total score:

Male time taken (n=311):
  Shapiro-Wilk test for samples between 200 and 2000 entities: W=0.9374,
p=0.0000
  Normal distribution: False

Male total score (n=311):
  Shapiro-Wilk test for samples between 200 and 2000 entities: W=0.9646,
p=0.0000
  Normal distribution: False

Both variables normally distributed: False

Testing normality for Female time taken and Female total score:

Female time taken (n=258):
  Shapiro-Wilk test for samples between 200 and 2000 entities: W=0.8720,
p=0.0000
  Normal distribution: False

Female total score (n=258):
  Shapiro-Wilk test for samples between 200 and 2000 entities: W=0.9762,
p=0.0003
  Normal distribution: False

Both variables normally distributed: False

Gender-specific results:
  Male (Spearman): ρ = -0.0455, p = 0.4235
  Female (Spearman): ρ = -0.1996, p = 0.0013
    \end{Verbatim}

    \subsubsection{5.2 Correlation Results
Interpretation}\label{correlation-results-interpretation}

The following interpretation explains the time-score relationship
findings in accessible language:

    \begin{Verbatim}[commandchars=\\\{\}]
UNDERSTANDING TIME-SCORE RELATIONSHIP RESULTS
======================================================================

 BASIC INFORMATION:
   • Average exam time: 200.5 minutes
   • Average exam score: 64.9 points
   • Total students analyzed: 569

 QUESTION 1: Is there a relationship between TIME SPENT and PERFORMANCE?
   ANSWER: There is a weak relationship
   Direction: Students who spend MORE time tend to score LOWER
   What this means: Time spent has a small connection to exam scores
   Reliability: This finding is very reliable (99\%+ confidence)

 QUESTION 2: What does this mean for EXAM STRATEGY?
   STRATEGY: Time has a small but real effect on performance
   • Taking more time is associated with lower performance (may indicate
struggling)

 QUESTION 3: Are there GENDER DIFFERENCES in the time-score relationship?
   ANSWER: The time-score relationship is somewhat different between genders
   Male correlation: -0.046 (p = 0.424)
   Female correlation: -0.200 (p = 0.001)
   What this means: Time has a stronger effect on female student performance

======================================================================
    \end{Verbatim}

    \subsubsection{5.3 Time Analysis
Visualizations}\label{time-analysis-visualizations}

The following visualizations examine time usage patterns by gender and
their relationship with exam performance.

    \paragraph{5.3.1 Time Statistics by
Gender}\label{time-statistics-by-gender}

    \begin{Verbatim}[commandchars=\\\{\}]
=== TIME ANALYSIS DESCRIPTIVE STATISTICS ===
             Male   Female
Count     311.000  258.000
Mean      190.500  212.580
Median    198.780  228.690
Mode      240.030  240.070
Std Dev    47.870   37.750
Min        40.750   81.930
Max       278.280  297.530
Q1        159.680  195.760
Q3        234.360  239.120
IQR        74.680   43.360
Skewness   -0.680   -1.106
Kurtosis   -0.246    0.700
    \end{Verbatim}

    \paragraph{5.3.2 Time Distribution and Time-Score Relationship
Plots}\label{time-distribution-and-time-score-relationship-plots}

The following visualizations provide different perspectives on time
usage: 1. \textbf{Boxplot}: Time distribution by gender with outlier
identification 2. \textbf{Scatter plot}: Time vs score relationship by
gender

    \begin{Verbatim}[commandchars=\\\{\}]
=== TIME OUTLIER ANALYSIS ===
Male Time: 1 outliers (0.3\%)
  Range: 40.75 - 40.75
Female Time: 9 outliers (3.5\%)
  Range: 81.93 - 130.53
    \end{Verbatim}

    \begin{center}
    \adjustimage{max size={0.9\linewidth}{0.9\paperheight}}{3b. Overall_exam_performance_analysis_genders_files/3b. Overall_exam_performance_analysis_genders_37_1.png}
    \end{center}
    { \hspace*{\fill} \\}
    
    
    \begin{Verbatim}[commandchars=\\\{\}]
<Figure size 640x480 with 0 Axes>
    \end{Verbatim}

    

    % Add a bibliography block to the postdoc
    
    
    
\end{document}
